\section{Order Types and Matching}
\label{sec:order-types}

\Zoo{} \DEX{} ships the full retail and institutional order-type
matrix on day one of the April 20, 2026 launch.

\subsection{Order types (taker side)}

\paragraph{Limit.}
Buy or sell at a specified price or better. Resting limit orders
post to the continuous limit book at the specified price level and
match by price-time priority.

\paragraph{Market.}
Buy or sell at the best available price. Routes against the AMM
pool when the AMM has a better price than the top of the book; uses
the order book otherwise.

\paragraph{Stop.}
Triggered when the market reaches the stop price; converts into a
market or limit order at the trigger.

\paragraph{IOC (Immediate-or-Cancel).}
Match what can be matched at the specified price; cancel any
unfilled remainder. No resting state.

\paragraph{FOK (Fill-or-Kill).}
Match the full quantity at the specified price or cancel the entire
order. No partial fills.

\paragraph{Post-only.}
Add liquidity to the order book at the specified price; reject if
the order would cross the spread (i.e., act as a taker). Useful for
makers earning rebates.

\paragraph{Hidden.}
Resting limit order that does not display in the public order book.
Visible to the matching engine only. Used by institutional flow that
does not want to telegraph size.

\paragraph{Iceberg.}
A large limit order that displays only a fraction of its size in the
public book. As the displayed slice fills, the next slice becomes
visible. Useful for working size without spooking the book.

\subsection{Continuous limit book}

The continuous limit book is the high-frequency matching engine. It
implements price-time priority: among orders at the same price, the
earliest-arrived order matches first. The order book lives in device
memory per LP-137 \cite{lp-137}; updates run as \GPU{} kernels.

A typical continuous-book trade lifecycle:

\begin{enumerate}[leftmargin=1.4em]
  \item Order arrives at the sequencer; passes the \Liquidity{}
        precompile check.
  \item Order is inserted into the per-pair price-time-priority
        order book.
  \item If the order crosses an existing book level, the crossing
        match is generated; both orders' resting states are updated;
        the trade event is emitted.
  \item Trade event flows through the transfer-agent notify path
        \cite{liquidity-protocol-formal}.
  \item Block production seals the trades into a \Quasar{} 4.0
        triple-cert.
\end{enumerate}

\subsection{Batch auctions}

For fair price discovery on instruments where front-running risk is
material (illiquid securities, opening prints, closing prints), the
\DEX{} also runs batch auctions. A batch auction:

\begin{enumerate}[leftmargin=1.4em]
  \item Accepts orders during the batch window (e.g., 100ms).
  \item At the window close, computes the uniform clearing price
        that maximises matched volume.
  \item Executes all crossing orders at the uniform clearing price.
  \item Cancels any orders that do not match at the clearing price
        (or rests them in the continuous book if so flagged).
\end{enumerate}

Batch auctions eliminate the latency arms race because the inside
order does not benefit from being earlier within the batch; only
prices and quantities matter, not microsecond arrival times. This
is materially better for retail than continuous matching for the
same instrument set.

\subsection{When to use which}

\begin{longtable}{p{0.30\textwidth}p{0.62\textwidth}}
\toprule
\textbf{Mode} & \textbf{Use when} \\
\midrule
Continuous limit book & Liquid pairs, high-frequency strategies, market-makers, retail at-the-touch trading. \\
Batch auction & Opening prints, closing prints, illiquid securities, retail order flow that benefits from front-running protection. \\
\AMM{} (V2/V3) & Long-tail pairs, programmatic-LP yield, automated market-making without an active maker. \\
\bottomrule
\end{longtable}

The matching engine routes a market order across all three venues
and executes against whichever has the best price; the user receives
a single fill report.

\subsection{Self-trade prevention}

Two orders from the same beneficial owner cannot match each other.
The matching engine consults the \Liquidity{} precompile's
beneficial-owner record; if a match would put both sides on the same
beneficial owner, the engine cancels the resting side and applies
the inbound side as a new resting order. This prevents wash trading
without requiring the user to explicitly tag orders.

\subsection{Order routing across V2/V3/V4}

A retail market order for a security pair flows through the router:

\begin{enumerate}[leftmargin=1.4em]
  \item Router queries V2 pool, V3 pool, V4 (Lux \DEX{}), and the
        local CLOB precompile for executable price at the order
        size.
  \item Router computes the price-impact-adjusted best venue (or
        splits across venues for large orders).
  \item Router executes against the chosen venue(s); each leg goes
        through the \Liquidity{} precompile check.
  \item Router returns a single weighted-average fill report to the
        user.
\end{enumerate}

This is invisible to the user; it is a router-side concern.

\subsection{Cancel-replace semantics}

A cancel-replace transaction atomically cancels the existing order
and submits a new one. If the cancel succeeds and the replace fails
(e.g., the precompile rejects the new order), the engine reverts
to a clean state with no resting order. There is no half-applied
state.

\subsection{Reporting}

Every executed trade emits a \texttt{TradeEvent} on chain that
includes:

\begin{itemize}[leftmargin=1.2em]
  \item Pair, price, quantity, side, taker/maker flag.
  \item Beneficial-owner reference (or its hash, in privacy mode).
  \item Venue (V2 / V3 / V4 / CLOB).
  \item \Liquidity{}-precompile attestation hash.
  \item \Quasar{}Cert reference (for finality verification).
\end{itemize}

Public market-data subscribers (the BBO, the last-trade tape, the
volume by venue) consume these events. There is no separate
market-data infrastructure; the chain log is the tape.
